\documentclass[11pt]{article}
\usepackage[a4paper,margin=28mm]{geometry}
\usepackage{amsmath,amssymb,amsthm}
\usepackage[colorlinks=true,linkcolor=blue,urlcolor=blue,citecolor=blue]{hyperref}
\newtheorem{theorem}{Theorem}
\newcommand{\R}{\mathbb R}
\newcommand{\C}{\mathbb C}
\newcommand{\Q}{\mathbb Q}
\newcommand{\Z}{\mathbb Z}
\newcommand{\F}{\mathbb F}
\DeclareMathOperator{\End}{End}
\DeclareMathOperator{\Hom}{Hom}
\DeclareMathOperator{\Aut}{Aut}
\DeclareMathOperator{\Gal}{Gal}
\DeclareMathOperator{\Ad}{Ad}
\DeclareMathOperator{\GL}{GL}
\DeclareMathOperator{\PGL}{PGL}
\title{Excluding torsion of orders \(26\) and \(39\) on geometrically split abelian surfaces}
\author{}
\date{21 October 2026}


\begin{document}
\maketitle

\begin{theorem}\label{thm:main}
Let \(A/\Q\) be a geometrically split abelian surface. Then \(A(\Q)\)
contains no point of order \(26\) and no point of order \(39\).
In particular, none of
\href{https://genus-2-torsion.github.io/group.html?g=26}{\([26]\)},
\href{https://genus-2-torsion.github.io/group.html?g=2.26}{\([2,26]\)} and
\href{https://genus-2-torsion.github.io/group.html?g=39}{\([39]\)}
occurs as the rational torsion subgroup of a geometrically split
genus-\(2\) Jacobian over \(\Q\).
\end{theorem}

Here ``geometrically split'' means that \(A_{\overline{\Q}}\) is
isogenous to a product of two elliptic curves. There is no restriction
on the degree of this isogeny, on the field of definition of the
geometric endomorphisms, or on the integral endomorphism ring, and the
theorem does not require a principal polarization. Along the way we
prove two statements about non-CM geometric squares whose endomorphism
field is cyclic of degree \(6\) (Proposition~9.1 and Lemma~10.1); they
need no torsion hypothesis. One step uses a Magma computation on the
modular curve \(X_0(39)\):
\href{https://genus-2-torsion.github.io/data/knowledge/sources/geometric_26_39_x039_w3.m}{\texttt{geometric\_26\_39\_x039\_w3.m}}
with its
\href{https://genus-2-torsion.github.io/data/knowledge/sources/geometric_26_39_x039_w3.log}{log}.
The script
\href{https://genus-2-torsion.github.io/data/knowledge/sources/geometric_26_39_checks.m}{\texttt{geometric\_26\_39\_checks.m}}
(\href{https://genus-2-torsion.github.io/data/knowledge/sources/geometric_26_39_checks.log}{log})
contains supplementary checks for Sections~4 and~11; the proof does
not depend on it.

\section{Inputs}

We use the following results on torsion points of elliptic curves.
\begin{enumerate}
\item Mazur \cite[Theorem~(7$'$)]{Mazur}: if an elliptic curve over
\(\Q\) has a rational point of order \(m\), then \(m\le10\) or
\(m=12\). In particular, no elliptic curve over \(\Q\) has a rational
point of order \(13\).
\item Kenku--Momose \cite[Theorem~(0.1)]{KM}: let \(N\) be a product of
powers of \(2,3,5,7,11,13\) and let \(m\mid N\). If \(X_1(m,N)\) has
positive genus and \((m,N)\) is not one of \((1,11)\), \((1,13)\),
\((1,14)\), \((1,15)\), \((1,16)\), \((1,18)\), \((2,10)\), \((2,12)\),
then \(Y_1(m,N)(k)=\emptyset\) for every quadratic field \(k\). Taking
\((m,N)=(1,26)\) and \((1,39)\) (the curves \(X_1(26)\) and
\(X_1(39)\) have genus \(10\) and \(33\)), no elliptic curve over a
quadratic field has a rational point of order \(26\) or \(39\).
\item Derickx--Etropolski--van Hoeij--Morrow--Zureick-Brown
\cite[Theorem~A]{DEvHMZB}: if \(E\) is an elliptic curve over a cubic
field \(K\), then \(E(K)_{\mathrm{tors}}\) is \(\Z/N\) with
\(N\in\{1,\ldots,16,18,20,21\}\), or \(\Z/2\times\Z/2N\) with
\(1\le N\le7\). None of these groups has an element of order \(26\) or
\(39\).
\item Derickx--Najman \cite[Theorem~1.1]{DN}: over quartic fields the
torsion groups of elliptic curves are \(\Z/n\) with
\(n=1,\ldots,18,20,21,22,24\), and \(\Z/2\times\Z/2n\) \((n\le9)\),
\(\Z/3\times\Z/3n\) \((n\le3)\), \(\Z/4\times\Z/4n\) \((n\le2)\),
\(\Z/5\times\Z/5\), \(\Z/6\times\Z/6\). None of them has an element of
order \(26\) or \(39\).
\end{enumerate}
Together, these give:
\[
\text{if } [K:\Q]\le 4, \text{ no elliptic curve over } K
\text{ has a } K\text{-rational point of order } 26 \text{ or } 39.
\tag{1.1}
\]

We also use the following.
\begin{itemize}
\item Faltings \cite[\S5, Satz~3, Satz~4 and Korollar~1]{Faltings}: for
abelian varieties \(X,Y\) over a number field \(K\) and a prime
\(\ell\), the representation of \(G_K=\Gal(\overline K/K)\) on
\(V_\ell(X)=T_\ell(X)\otimes\Q_\ell\) is semisimple, and the natural
map
\[
\Hom_K(X,Y)\otimes\Z_\ell\longrightarrow\Hom_{G_K}(T_\ell X,T_\ell Y)
\]
is an isomorphism.
\item The criterion of N\'eron--Ogg--Shafarevich
\cite[\S1, Theorem~1]{SerreTate}: an abelian variety over a finite
extension \(F\) of \(\Q_p\) has good reduction if and only if the
inertia group \(I_F\) acts trivially on \(T_\ell\) for one
(equivalently, every) prime \(\ell\ne p\). Moreover, an elliptic curve
with complex multiplication over a number field has potentially good
reduction at every prime \cite[\S5, Theorem~6(a)]{SerreTate}.
\item Weil's theorem: if \(\bar A\) is an abelian surface over
\(\F_q\), then \(\#\bar A(\F_q)=f(1)\), where \(f\in\Z[T]\) is the
characteristic polynomial of Frobenius; it is monic of degree \(4\),
and all its complex roots have absolute value \(\sqrt q\).
\item Lattice independence (Brauer--Nesbitt): if a compact group acts
continuously on a finite-dimensional \(\Q_p\)-vector space, the
reductions modulo \(p\) of any two stable lattices have isomorphic
semisimplifications. This applies to isomorphic representations after
transporting a lattice through the isomorphism. We also use the
Skolem--Noether theorem and Clifford's theorem; the latter is proved
where it is used.
\end{itemize}

\section{Four types and two elementary lemmas}

For a prime \(p\) let \(T_p(A)\) be the Tate module and
\(V_p(A)=T_p(A)\otimes\Q_p\), with their actions of
\(G_{\Q}=\Gal(\overline\Q/\Q)\). For a stable lattice
\(T'\subset V_p(A)\), the module \(T'/pT'\) is its residual
representation, and \((\cdot)^{\mathrm{ss}}\) denotes
semisimplification. By lattice independence, \(A[p]^{\mathrm{ss}}\)
depends only on the \(\Q\)-isogeny class of \(A\). A rational point of
order \(p\) spans a trivial submodule of \(A[p]=T_p(A)/pT_p(A)\); hence
the residual representation of every \(G_{\Q}\)-stable lattice in
\(V_p(A)\) has a trivial composition factor. We write \(\bar\chi_p\)
for the mod-\(p\) cyclotomic character.

A geometrically split abelian surface \(A/\Q\) is of exactly one of
the following types.
\begin{itemize}
\item[\textbf{(Q)}] \(A\) is not simple over \(\Q\).
\item[\textbf{(N)}] \(A\) is simple over \(\Q\), and
\(A_{\overline\Q}\sim E_1\times E_2\) with \(E_1,E_2\) not isogenous.
\item[\textbf{(CM)}] \(A\) is simple over \(\Q\), and
\(A_{\overline\Q}\sim E^2\) with \(E\) having complex multiplication.
\item[\textbf{(S)}] \(A\) is simple over \(\Q\), and
\(A_{\overline\Q}\sim E^2\) with \(E\) without complex multiplication.
\end{itemize}

Suppose from now on that \(P\in A(\Q)\) has order \(N\in\{26,39\}\).
Write \(N=13s\) with \(s\in\{2,3\}\), and put \(P_{13}=sP\) and
\(P_s=13P\); these are rational points of orders \(13\) and \(s\).

\medskip\noindent\textbf{Lemma 2.1 (coprime torsion).}
\emph{Let \(K\) be a number field and \(m,n\) coprime positive
integers. Let \(E_m,E_n\) be \(K\)-isogenous elliptic curves over \(K\)
such that \(E_m(K)\) has a point of order \(m\) and \(E_n(K)\) has a
point of order \(n\). Then some elliptic curve over \(K\) that is
\(K\)-isogenous to \(E_m\) has a \(K\)-rational point of order \(mn\).}

\begin{proof}
Let \(f:E_m\to E_n\) be a \(K\)-isogeny, and let \(\Phi\subset\ker f\)
be the subgroup of elements whose order is divisible only by primes
dividing \(n\). It is characteristic in \(\ker f\), hence
\(G_K\)-stable. Factor \(f\) as
\[
E_m\xrightarrow{\ a\ }E'=E_m/\Phi\xrightarrow{\ b\ }E_n .
\]
The degree of \(a\) is prime to \(m\), so \(a\) maps the given point
of order \(m\) to a \(K\)-point of order \(m\) on \(E'\). The degree
of \(b\) is prime to \(n\). Since \(b\circ\hat b=[\deg b]\), the dual
isogeny \(\hat b\) is injective on \(E_n[n]\), so it maps the given
point of order \(n\) to a \(K\)-point of order \(n\) on \(E'\). The
sum of the two points has order \(mn\).
\end{proof}

\noindent\textbf{Lemma 2.2.}
\emph{Let \(k\) be a field of characteristic \(0\), let \(X\) be an
abelian surface over \(k\), let \(E\subset X\) be an elliptic curve
defined over \(k\), and put \(F=X/E\). If \(Q\in X(k)\) has prime
order \(q\), then either \(Q\in E(k)\), or the image of \(Q\) in
\(F(k)\) has order \(q\).}

\begin{proof}
The image of \(Q\) is killed by \(q\), and it is zero exactly when
\(Q\in E\).
\end{proof}

\section{Surfaces that are not simple over \texorpdfstring{\(\Q\)}{Q}, and nonisogenous factors}

\textbf{Type (Q).} If \(A\) is not simple over \(\Q\), it contains an
elliptic curve \(E\) defined over \(\Q\), and \(A/E\) is an elliptic
curve over \(\Q\). By Lemma~2.2 with \(k=\Q\), the point \(P_{13}\)
gives a rational point of order \(13\) on \(E\) or on \(A/E\). This
contradicts Mazur's theorem.

\textbf{Type (N).} Suppose \(A_{\overline\Q}\sim E_1\times E_2\) with
\(E_1\not\sim E_2\). For each \(i\), all nonzero homomorphisms
\(E_i\to A_{\overline\Q}\) have the same image:
\(\Hom(E_i,A_{\overline\Q})\otimes\Q\) is a vector space of dimension
one over \(\End^0(E_i)\), because \(\Hom(E_1,E_2)=0\). Every elliptic
curve in \(A_{\overline\Q}\) is isogenous to \(E_1\) or to \(E_2\).
Hence \(A_{\overline\Q}\) has exactly two elliptic subvarieties. The
group \(G_{\Q}\) permutes them. If it fixed them, they would be defined
over \(\Q\), and \(A\) would not be simple over \(\Q\). Hence they are
exchanged, and they are defined over a quadratic field \(K\). Write
them as \(E\) and \(E^\sigma\), where \(\sigma\) generates
\(\Gal(K/\Q)\).

Put \(F=A_K/E\). Since \(E\cap E^\sigma\) is finite and
\(E+E^\sigma=A\), the quotient map restricts to a \(K\)-isogeny
\(E^\sigma\to F\); hence \(F^\sigma\) is \(K\)-isogenous to \(E\).
Apply Lemma~2.2 over \(K\) to \(P_{13}\) and to \(P_s\). If the point
lands in \(F(K)\), applying \(\sigma\) gives a point of the same order
in \(F^\sigma(K)\). Thus the \(K\)-isogeny class of \(E\) contains
curves with \(K\)-points of orders \(13\) and \(s\). Lemma~2.1 gives an
elliptic curve over the quadratic field \(K\) with a \(K\)-point of
order \(N\), contradicting (1.1).

\section{The CM case: a local argument}

\noindent\textbf{Lemma 4.1.}
\emph{Let \(p\) be a prime and \(Y/\Q_p\) an abelian variety with
potentially good reduction. There is an abelian variety \(\bar Y\) over
\(\F_p\), of the same dimension, such that \(Y(\Q_p)[M]\) embeds into
\(\bar Y(\F_p)\) for every integer \(M\) prime to \(p\).}

\begin{proof}
This is the argument of Clark \cite[\S2.1]{Clark} and of
Clark--Xarles \cite[Main Theorem~(i) and \S2.4]{ClarkXarles}, who quote
the first step from \cite[p.~498]{SerreTate}. We include it for
completeness.

\emph{A totally ramified extension.} Fix a prime \(\ell\ne p\), let
\(\rho:G_{\Q_p}\to\GL(T_\ell Y)\) be the Galois representation, and
let \(I\subset G_{\Q_p}\) be the inertia group. Let \(F'/\Q_p\) be a
finite extension over which \(Y\) has good reduction. By the
N\'eron--Ogg--Shafarevich criterion, \(I\cap G_{F'}\) acts trivially
on \(T_\ell Y\). It has finite index in \(I\), so \(\rho(I)\) is
finite. The group \(N_0=I\cap\ker\rho\) is closed and normal in
\(G_{\Q_p}\). Put \(\mathcal G=G_{\Q_p}/N_0\). Then \(I/N_0\) is a
finite normal subgroup of \(\mathcal G\), and
\(\mathcal G/(I/N_0)\simeq G_{\Q_p}/I\simeq\widehat\Z\). Since
\(\widehat\Z\) is a free profinite group, the projection
\(\mathcal G\to\widehat\Z\) has a continuous section. Let \(S\) be its
image. Then \(S\cap(I/N_0)=1\) and \(S\cdot(I/N_0)=\mathcal G\). Let
\(F\) be the fixed field of the preimage of \(S\) in \(G_{\Q_p}\).
Then \([F:\Q_p]=\#(I/N_0)\) is finite. The field
\(F\cap\Q_p^{\mathrm{nr}}\) is fixed by \(G_F\) and by \(I\), whose
images generate \(\mathcal G\). Hence \(F\cap\Q_p^{\mathrm{nr}}=\Q_p\),
so \(F/\Q_p\) is totally ramified and has residue field \(\F_p\). The
inertia group \(G_F\cap I\) maps into \(S\cap(I/N_0)=1\), so it lies in
\(\ker\rho\). By the N\'eron--Ogg--Shafarevich criterion, \(Y_F\) has
good reduction.

\emph{Injectivity.} Let \(\mathcal Y\) be the abelian scheme over the
ring of integers \(\mathcal O_F\) with generic fibre \(Y_F\), and let
\(\bar Y\) be its special fibre, an abelian variety over \(\F_p\). For
\(M\) prime to \(p\), the group scheme \(\mathcal Y[M]\) is finite
\'etale over the henselian local ring \(\mathcal O_F\). Hence reduction
\(\mathcal Y[M](\mathcal O_F)\to\bar Y[M](\F_p)\) is injective. By
properness, \(Y(F)[M]=\mathcal Y[M](\mathcal O_F)\). Thus
\(Y(\Q_p)[M]\subset Y(F)[M]\) embeds into \(\bar Y(\F_p)\).
\end{proof}

\noindent\textbf{Proposition 4.2.}
\emph{If \(A_{\overline\Q}\) is isogenous to \(E^2\) for an elliptic
curve \(E\) with complex multiplication, then \(A(\Q)\) has no point of
order \(26\) or \(39\).}

\begin{proof}
Let \(p\) be a prime. Let \(M\) be a number field over which \(E\) and
an isogeny \(A_M\to E_M^2\) are defined, and let \(v\) be a place of
\(M\) above \(p\). Since \(E\) has potentially good reduction at \(v\),
there is a finite extension \(F'\) of \(M_v\) over which \(E\) has good
reduction. Then \(V_\ell(A)\simeq V_\ell(E)^2\) as representations of
\(G_{F'}\) for \(\ell\ne p\), so the inertia group of \(F'\) acts
trivially on \(T_\ell(A)\). By the N\'eron--Ogg--Shafarevich criterion,
\(A\) has good reduction over \(F'\). Thus \(A_{\Q_p}\) has potentially
good reduction, and Lemma~4.1 applies to it, giving an abelian surface
\(\bar A\) over \(\F_p\).

\emph{Order 39.} Take \(p=2\). Since \(39\) is odd, \(\Z/39\) embeds
into \(\bar A(\F_2)\), so \(\#\bar A(\F_2)\ge39\). But Weil's theorem
gives \(\#\bar A(\F_2)\le(1+\sqrt2)^4=17+12\sqrt2<34\), a
contradiction.

\emph{Order 26.} Take \(p=3\). Then \(\Z/26\) embeds into
\(\bar A(\F_3)\), so \(26\) divides \(\#\bar A(\F_3)\), and
\(\#\bar A(\F_3)\le(1+\sqrt3)^4=28+16\sqrt3<56\). Hence
\(\#\bar A(\F_3)\in\{26,52\}\). Write the characteristic polynomial of
Frobenius as
\[
f(T)=T^4+aT^3+bT^2+3aT+9,\qquad a,b\in\Z .
\]
Then \(f(T)=T^2h(T+3/T)\) with \(h(X)=X^2+aX+(b-6)\). If \(c\) is a
root of \(h\), the roots \(\pi\) of \(T^2-cT+3\) are roots of \(f\),
so \(|\pi|=\sqrt3\) and \(c=\pi+3/\pi=\pi+\bar\pi\) is real with
\(|c|\le2\sqrt3\). Let \(c_1,c_2\) be the roots of \(h\). Then
\(a=-(c_1+c_2)\) and \(b=6+c_1c_2\). The bounds \(|c_i|\le2\sqrt3\),
the inequality \(c_1c_2\le\big((c_1+c_2)/2\big)^2\) and the
inequalities \((2\sqrt3\mp c_1)(2\sqrt3\mp c_2)\ge0\) give
\[
|a|\le4\sqrt3,\qquad 2\sqrt3\,|a|-6\ \le\ b\ \le\ \frac{a^2}{4}+6.
\tag{4.1}
\]
We have \(\#\bar A(\F_3)=f(1)=10+4a+b\).
\begin{itemize}
\item If \(f(1)=26\), then \(b=16-4a\). The upper bound in (4.1) gives
\(a^2+16a-40\ge0\). Together with \(|a|\le4\sqrt3\) this forces
\(a\ge-8+\sqrt{104}\), so \(a\ge3\). The lower bound then reads
\((4+2\sqrt3)\,a\le22\), so \(a\le2\). This is a contradiction.
\item If \(f(1)=52\), then \(b=42-4a\). The upper bound gives
\(a^2+16a-144\ge0\). With \(|a|\le4\sqrt3\) this forces
\(a\ge-8+\sqrt{208}>6\), so \(a\ge7>4\sqrt3\). This is a
contradiction.\qedhere
\end{itemize}
\end{proof}

The supplementary script confirms these two computations by listing
every value of \(f(1)\) allowed by the condition on the roots of \(h\):
over \(\F_3\) the values \(26\) and \(52\) do not occur, and over
\(\F_2\) all values are at most \(25\). As an independent check, Clark
\cite[Theorem~5]{Clark} (the same list is
\cite[Theorem~1.4]{ClarkXarles}) shows that the torsion order of an
abelian surface over \(\Q\) with everywhere potentially good reduction
lies in
\[
\{1,\ldots,16,18,19,20,22,24,25,28,30,36,48,60,72\};
\]
no element of this set is divisible by \(26\) or \(39\).

\section{The non-CM square case: the endomorphism field}

From now on, \(A\) has type (S). Put
\[
D=\End^0(A_{\overline\Q})=\End(A_{\overline\Q})\otimes\Q\simeq M_2(\Q).
\]
Let \(L\) be the smallest field over which all geometric endomorphisms
of \(A\) are defined. It is a finite Galois extension of \(\Q\). Put
\(\Gamma=\Gal(L/\Q)\). The group \(G_{\Q}\) acts on \(D\) through
\(\Gamma\), and \(\Gamma\) acts faithfully. By Skolem--Noether, each
\(g\in\Gamma\) acts as \(\Ad(R_g)\) for some \(R_g\in\GL_2(\Q)\),
unique up to \(\Q^\times\). This embeds \(\Gamma\) into \(\PGL_2(\Q)\);
we write \(g=[R_g]\). Note that \(\End^0_{\Q}(A)=D^{\Gamma}\), and that
\(\Gamma\ne1\): otherwise every elliptic subvariety of
\(A_{\overline\Q}\) would be defined over \(\Q\) (Lemma~5.2 below),
contradicting simplicity over \(\Q\).

\medskip\noindent\textbf{Lemma 5.1 (finite subgroups of
\(\PGL_2(\Q)\)).}
\emph{\(\Gamma\) is one of \(C_1,C_2,C_3,C_4,C_6\),
\(D_4\simeq C_2\times C_2\), \(D_6,D_8,D_{12}\), where \(C_n\) is
cyclic of order \(n\) and the subscript of \(D\) is the group order. If
\([R]\in\PGL_2(\Q)\) has order \(n\in\{3,4,6\}\), then
\(\operatorname{tr}R\ne0\), and the unique representative with trace
\(1\) has characteristic polynomial \(X^2-X+1/c\), where \(c=1,2,3\)
for \(n=3,4,6\); its discriminant is \(-3,-1,-1/3\), respectively. If
\([R]\) has order \(2\), then \(\operatorname{tr}R=0\) and
\(R^2\in\Q^\times\).}

\begin{proof}
A finite subgroup of \(\PGL_2(\R)\) fixes a point of the hyperbolic
plane, and its action on the tangent plane identifies it with a finite
subgroup of \(O(2)\). Hence it is cyclic or dihedral. If \([R]\) has
order \(n>2\), the ratio of the eigenvalues of \(R\) is a primitive
\(n\)-th root of unity \(\zeta\), and
\[
\frac{(\operatorname{tr}R)^2}{\det R}=2+\zeta+\zeta^{-1}\in\Q,
\]
which forces \(n\in\{3,4,6\}\), with value \(c=1,2,3\). Normalizing
\(\operatorname{tr}R=1\) gives \(\det R=1/c\). If \([R]\) has order
\(2\), then \(R^2\) is a scalar while \(R\) is not, so the
characteristic polynomial of \(R\) is \(X^2-R^2\).
\end{proof}

In a dihedral group \(D_{2n}\) with \(n\in\{3,4,6\}\), let \(r\)
generate the rotation subgroup and let \(R\) be its trace-one
representative. A reflection \(t\) satisfies \(trt^{-1}=r^{-1}\). Since
\(R^{-1}=(1-R)/\det R\), the trace-one representative of \(r^{-1}\) is
\(1-R\). Conjugation preserves the trace, so \(\Ad(R_t)(R)=1-R\).

\medskip\noindent\textbf{Lemma 5.2.}
\emph{Every elliptic curve \(E\subset A_{\overline\Q}\) is defined over
\(L\), and \(A_L/E\) is \(L\)-isogenous to \(E\).}

\begin{proof}
The curve \(E\) is the image of \(n\epsilon\) for a rank-one idempotent
\(\epsilon\in D\) and a suitable integer \(n\ge1\). This endomorphism
is defined over \(L\), so \(E\) is defined over \(L\). Let \(E'\) be
the image of \(n(1-\epsilon)\). Then \(E+E'=A\) and \(E\cap E'\) is
finite, so the quotient map restricts to an \(L\)-isogeny
\(E'\to A_L/E\). Since \(D=M_2(\Q)\), we have
\(\epsilon D(1-\epsilon)\ne0\), so there is a nonzero homomorphism
between \(E'\) and \(E\). It is defined over \(L\). Hence
\(E\sim_L E'\sim_L A_L/E\).
\end{proof}

\section{Three lemmas on residual representations}

In this section \(A\) has type (S), \(p\) is an odd prime,
\(T=T_p(A)\) and \(V=V_p(A)\). The algebra
\(D_p=D\otimes\Q_p\simeq M_2(\Q_p)\) acts on \(V\), and for
\(g\in G_{\Q}\), \(\varphi\in D_p\), \(v\in V\) we have
\(g(\varphi v)=g(\varphi)\,g(v)\). As a module over \(M_2(\Q_p)\),
\(V\) is the sum of two copies of the standard module \(\Q_p^2\).

\medskip\noindent\textbf{Lemma 6.1 (polarization).}
\emph{Let \(T'\subset V\) be a \(G_{\Q}\)-stable lattice such that
every composition factor of \((T'/pT')|_{G_L}\) is trivial. Then
\(\Q(\zeta_p)\subset L\).}

\begin{proof}
Choose a polarization of \(A\) defined over \(\Q\); every abelian
variety over \(\Q\) has one. Its Weil pairing is a nondegenerate
alternating \(G_{\Q}\)-equivariant pairing \(V\times V\to\Q_p(1)\), so
\(V\simeq V^*\otimes\Q_p(1)\). Under this isomorphism the lattice
\(\Hom(T',\Z_p(1))\) becomes a \(G_{\Q}\)-stable lattice in \(V\)
whose reduction is \((T'/pT')^*\otimes\bar\chi_p\). Restricted to
\(G_L\), its composition factors are four copies of
\(\bar\chi_p|_{G_L}\). Lattice independence over \(G_L\) gives
\(\bar\chi_p|_{G_L}=1\), that is, \(\Q(\zeta_p)\subset L\). The degree
of the polarization plays no role.
\end{proof}

\noindent\textbf{Lemma 6.2 (unramified order).}
\emph{Let \(\mathcal B\subset D_p\) be a \(\Z_p\)-order isomorphic to
the ring of integers of the unramified quadratic extension of \(\Q_p\).
Suppose that \(\Ad(R_g)(\mathcal B)=\mathcal B\) for every
\(g\in\Gamma\), and that some \(g_0\in\Gamma\) is represented by a unit
\(R_{g_0}\in\mathcal B^\times\) whose reduction lies in
\(\mathcal B/p\mathcal B\setminus\F_p\simeq\F_{p^2}\setminus\F_p\). If
\(A(\Q)[p]\ne0\), then \(\Q(\zeta_p)\subset L\).}

\begin{proof}
\begin{enumerate}
\item \emph{A stable maximal order.} Identify \(D_p\) with
\(\End_{\Q_p}(\Q_p^2)\). The \(\mathcal B\)-stable lattices in
\(\Q_p^2\) are the fractional ideals of the discrete valuation ring
\(\mathcal B\) times a fixed vector, so they form a single homothety
class \([\Lambda]\). For \(g\in\Gamma\), the lattice \(R_g\Lambda\) is
stable under \(R_g\mathcal BR_g^{-1}=\mathcal B\), so it lies in
\([\Lambda]\). Hence the maximal order
\(\mathcal R=\End_{\Z_p}(\Lambda)\simeq M_2(\Z_p)\) is
\(\Gamma\)-stable, and we may rescale each \(R_g\) so that
\(R_g\in\GL(\Lambda)\).
\item \emph{A tensor decomposition.} Put \(T'=\mathcal R\,T\). It is a
lattice, and it is \(G_{\Q}\)-stable because \(g(xt)=g(x)g(t)\) with
\(g(x)\in\mathcal R\). By Morita theory,
\(T'\simeq\Lambda\otimes_{\Z_p}W_0\), with \(\mathcal R\) acting on the
first factor and \(W_0=e_{11}T'\) free of rank \(2\).
\item \emph{Linear actions on the factors.} Let \(\widetilde G\) be the
group of pairs \((g,M)\) with \(g\in G_{\Q}\), \(M\in\GL(\Lambda)\) and
\(\Ad(M)=g\) on \(\mathcal R\). It surjects onto \(G_{\Q}\) with kernel
\(\Z_p^\times\). For \((g,M)\in\widetilde G\), the operator
\((M\otimes1)^{-1}g\) on \(T'\) commutes with \(\mathcal R\), so it
equals \(1\otimes n\) for a unique \(n\in\GL(W_0)\). Thus
\(\widetilde G\) acts linearly on \(\Lambda\) and on \(W_0\), and the
tensor product of these actions is the Galois action. Reducing modulo
\(p\), we get \(\bar V:=T'/pT'\simeq U\otimes W\) with
\(U=\Lambda/p\Lambda\) and \(W=W_0/pW_0\). The group \(\widetilde G\)
acts on \(\bar V\) through \(G_{\Q}\), so its submodules and
composition factors are the same for \(\widetilde G\) and for
\(G_{\Q}\).
\item \emph{\(U\) is irreducible.} Let \(g\in G_{\Q}\) map to \(g_0\).
Any \(M\) with \((g,M)\in\widetilde G\) equals \(\lambda R_{g_0}\) with
\(\lambda\in\Z_p^\times\). The \(\F_p\)-vector space \(U\) is a
one-dimensional \(\F_{p^2}\)-vector space, and \(M\) acts on it as
multiplication by an element of \(\F_{p^2}\setminus\F_p\). This
multiplication has no eigenvector over \(\F_p\), so \(U\) has no stable
line.
\item \emph{\(W\) is irreducible.} By lattice independence, \(\bar V\)
has a trivial composition factor. If \(W\) had a stable line \(\ell\),
then \(U\otimes\ell\) and \(U\otimes(W/\ell)\) would be twists of \(U\)
by characters, hence irreducible of dimension \(2\), and \(\bar V\)
would have no one-dimensional composition factor.
\item \emph{\(G_L\) acts trivially on \(\bar V\).} For \(h\in G_L\) the
action of \(h\) on \(\mathcal R\) is trivial, so
\((h,1)\in\widetilde G\). The subgroup \(H=\{(h,1):h\in G_L\}\) is
normal in \(\widetilde G\) and acts trivially on \(U\). Clifford's
theorem shows that \(W|_H\) is semisimple: if \(S\subset W\) is a
simple \(H\)-submodule, the translates \(xS\) for
\(x\in\widetilde G\) are simple \(H\)-submodules, and their sum is a
nonzero \(\widetilde G\)-submodule, hence equal to \(W\). Since
\(\bar V|_H\simeq W|_H^{\oplus2}\) has a trivial composition factor,
\(W^H\ne0\). As \(H\) is normal, \(W^H\) is \(\widetilde G\)-stable,
so \(W^H=W\).
\item Lemma~6.1 now gives \(\Q(\zeta_p)\subset L\).\qedhere
\end{enumerate}
\end{proof}

\noindent\textbf{Lemma 6.3 (split blocks).}
\emph{Let \(\mathcal B=\Z_pe_1\oplus\Z_pe_2\subset D_p\) with
orthogonal idempotents \(e_1+e_2=1\) of rank one. Suppose that
\(\Ad(R_g)(\mathcal B)=\mathcal B\) for every \(g\in\Gamma\) and that
some element of \(\Gamma\) exchanges \(e_1\) and \(e_2\). Let
\(K\subset L\) be the fixed field of the index-two subgroup fixing
\(e_1\). The group \(G_K\) acts on the line \(e_1D_pe_2\) by a
character \(\eta\), which factors through \(\Gal(L/K)\). Suppose that
its reduction \(\bar\eta\) is nontrivial. If \(A(\Q)[p]\ne0\), then
\(\Q(\zeta_p)\subset L\).}

\begin{proof}
\begin{enumerate}
\item \emph{Two blocks.} Put \(T'=\mathcal BT\), a \(G_{\Q}\)-stable
lattice, and \(T_i=e_iT'\), so \(T'=T_1\oplus T_2\). Each \(T_i\) has
rank \(2\). The group \(G_K\) preserves each \(T_i\), and every element
of \(G_{\Q}\setminus G_K\) exchanges them.
\item \emph{The twist relation.} Put \(V_i=T_i\otimes\Q_p\), and let
\(X\in e_1D_pe_2\) be nonzero. Then \(X:V_2\to V_1\) is an
isomorphism, and \(gXg^{-1}=\eta(g)X\) for \(g\in G_K\). Hence
\(V_1\simeq\eta\otimes V_2\) as \(G_K\)-representations. With
\(M_i=T_i/pT_i\), lattice independence gives
\(M_1^{\mathrm{ss}}\simeq\bar\eta\otimes M_2^{\mathrm{ss}}\). This does
not require \(X\) to map \(T_2\) onto \(T_1\).
\item \emph{Both blocks have a trivial factor.} The rational point of
order \(p\) gives a trivial \(G_{\Q}\)-composition factor of
\(T'/pT'=M_1\oplus M_2\). Restricting to \(G_K\), one of \(M_1,M_2\)
has a trivial \(G_K\)-composition factor. An element of
\(G_{\Q}\setminus G_K\) carries \(M_1\) to \(M_2\) and conjugates the
\(G_K\)-action, so both have one.
\item \emph{All factors are trivial on \(G_L\).} Write
\(M_2^{\mathrm{ss}}=\mathbf1\oplus\alpha\). Then
\(M_1^{\mathrm{ss}}=\bar\eta\oplus\bar\eta\alpha\) contains
\(\mathbf1\), and \(\bar\eta\ne1\), so \(\alpha=\bar\eta^{-1}\). All
four composition factors \(\mathbf1,\bar\eta^{-1},\bar\eta,\mathbf1\)
are trivial on \(G_L\).
\item Lemma~6.1 gives \(\Q(\zeta_p)\subset L\).\qedhere
\end{enumerate}
\end{proof}

\section{Reduction to a cyclic sextic endomorphism field}

\noindent\textbf{Proposition 7.1.}
\emph{Let \(A\) be of type (S) with a rational point of order
\(N\in\{26,39\}\). Then \(\Gamma\simeq C_6\).}

\begin{proof}
We first use only \(P_{13}\), and show that \(\Gamma\) is cyclic. We
have \(\Gamma\ne C_1\). Take \(p=13\).

\emph{Dihedral \(\Gamma=D_{2n}\), \(n\in\{3,4,6\}\).} Let \(r\)
generate the rotations and let \(R\) be its trace-one representative.
By Lemma~5.1, its characteristic polynomial is \(X^2-X+1/c\), with
discriminant \(-3\), \(-1\) or \(-1/3\). Since
\(13\equiv1\pmod{12}\), each of these is a nonzero square modulo
\(13\): \(-3\equiv6^2\), \(-1\equiv5^2\), \(-1/3\equiv2^2\). By
Hensel's lemma \(R\) has eigenvalues \(\alpha,\beta\in\Z_{13}^\times\)
with \(\alpha-\beta\in\Z_{13}^\times\), so
\(\mathcal B=\Z_{13}[R]=\Z_{13}e_1\oplus\Z_{13}e_2\) with
\(e_1=(R-\beta)/(\alpha-\beta)\) and \(e_2=1-e_1\). Rotations fix
\(R\); reflections send \(R\) to \(1-R\), which has eigenvalue
\(\beta\) on \(e_1\) and \(\alpha\) on \(e_2\). So \(\mathcal B\) is
\(\Gamma\)-stable, rotations fix \(e_1,e_2\), and reflections exchange
them. Here \(K=L^{\langle r\rangle}\) and \(\eta(r)=\alpha/\beta\), a
primitive \(n\)-th root of unity. Since \(n\mid12\), reduction modulo
\(13\) is injective on \(n\)-th roots of unity, so \(\bar\eta\ne1\).
Lemma~6.3 gives \(\Q(\zeta_{13})\subset L\). Then \(\Gamma\) surjects
onto \(\Gal(\Q(\zeta_{13})/\Q)\simeq C_{12}\), which is impossible,
because the abelianization of \(D_{2n}\) has exponent \(2\).

\emph{\(\Gamma=D_4=C_2\times C_2\).} The map
\([M]\mapsto v_{13}(\det M)\bmod2\) is a well-defined homomorphism
\(\Gamma\to\Z/2\). Choose \(r\ne1\) in its kernel, and let \(R\) be a
trace-zero representative. Then \(R^2=a\in\Q^\times\) with
\(v_{13}(a)\) even; after rescaling, \(a\in\Z_{13}^\times\). As
\(\Gamma\) is abelian, \(\Ad(R_g)(R)=\pm R\) for all \(g\in\Gamma\)
(both sides square to \(a\)), so \(\mathcal B=\Z_{13}[R]\) is
\(\Gamma\)-stable.
\begin{itemize}
\item If \(a\bmod13\) is not a square, \(\mathcal B\) is the ring of
integers of the unramified quadratic extension of \(\Q_{13}\), and
\(R\in\mathcal B^\times\) has reduction \(\bar R\) with
\(\bar R^2=\bar a\), so \(\bar R\notin\F_{13}\). Lemma~6.2 applies with
\(g_0=r\).
\item If \(a\bmod13\) is a nonzero square,
\(\mathcal B=\Z_{13}e_1\oplus\Z_{13}e_2\) with
\(e_{1,2}=(1\pm R/\sqrt a)/2\). The sign character \(g\mapsto\pm1\),
defined by \(\Ad(R_g)(R)=\pm R\), is nontrivial: otherwise every
\(R_g\) would lie in the centralizer
\((\Q_{13}[R])^\times\simeq\Q_{13}^\times\times\Q_{13}^\times\), and
\(\Gamma\) would embed in
\((\Q_{13}^\times\times\Q_{13}^\times)/\Q_{13}^\times\simeq\Q_{13}^\times\),
whose finite subgroups are cyclic. Its kernel is \(\langle r\rangle\),
elements outside it exchange \(e_1\) and \(e_2\), and
\(\eta(r)=\sqrt a/(-\sqrt a)=-1\), so \(\bar\eta\ne1\). Lemma~6.3
applies.
\end{itemize}
In both cases \(\Q(\zeta_{13})\subset L\), which is impossible since
\([L:\Q]=4<12\). Hence \(\Gamma\in\{C_2,C_3,C_4,C_6\}\).

\emph{\(\Gamma\in\{C_2,C_3,C_4\}\).} Then \([L:\Q]\in\{2,3,4\}\).
Choose an elliptic curve \(E\subset A_{\overline\Q}\). By Lemma~5.2,
\(E\) is defined over \(L\) and \(F=A_L/E\) is \(L\)-isogenous to
\(E\). Lemma~2.2 over \(L\), applied to \(P_{13}\) and \(P_s\), gives
\(L\)-points of orders \(13\) and \(s\) on curves in the
\(L\)-isogeny class of \(E\). Lemma~2.1 gives an elliptic curve over
\(L\) with an \(L\)-point of order \(N\), contradicting (1.1). Hence
\(\Gamma\simeq C_6\).
\end{proof}

\section{The cyclic sextic case: preliminaries}

In Sections~8--10, \(A/\Q\) is any abelian surface of type (S) with
\(\Gamma=\langle r\rangle\simeq C_6\); no torsion is assumed. Let \(R\)
be the trace-one representative of \(r\), so that
\(R^2-R+\tfrac13=0\), and let \(K=L^{\langle r^2\rangle}\) be the
quadratic subfield of \(L\). For \(g\in G_{\Q}\) write
\(g|_L=r^{k(g)}\) with \(k(g)\in\Z/6\). Then \(g\) acts on \(D\) as
\(\Ad(R^{k(g)})\), and \(G_K=\{g:k(g)\text{ even}\}\).

\subsection{The field \texorpdfstring{\(L\)}{L} is totally real}

Fix an embedding \(\overline\Q\subset\C\), and let \(c\in G_{\Q}\) be
complex conjugation. Put \(H=H_1(A(\C),\R)\), a real vector space of
dimension \(4\). The algebra \(D_\R=D\otimes\R\simeq M_2(\R)\) acts
faithfully on \(H\), so \(H\simeq U_\R\otimes W_\R\), with \(D_\R\)
acting on \(U_\R=\R^2\) and \(\dim W_\R=2\).
\begin{enumerate}
\item The complex structure \(J_H\) commutes with \(D_\R\), because the
endomorphisms are holomorphic. Hence \(J_H=1\otimes J\) with
\(J\in\End(W_\R)\) and \(J^2=-1\).
\item Since \(A\) is defined over \(\Q\), complex conjugation induces
an involution \(c_*\) of \(H\) with \(c_*J_H=-J_Hc_*\), and
\(c_*\varphi c_*^{-1}=c(\varphi)\) for \(\varphi\in D\). Thus \(c_*\)
normalizes \(D_\R\otimes1\). By Skolem--Noether there is
\(M\in\GL(U_\R)\) with
\(c_*(\varphi\otimes1)c_*^{-1}=M\varphi M^{-1}\otimes1\). Then
\((M\otimes1)^{-1}c_*\) commutes with \(D_\R\otimes1\), so
\(c_*=M\otimes N\) with \(N\in\GL(W_\R)\).
\item From \(c_*J_H=-J_Hc_*\) we get \(NJ=-JN\). In a basis with
\(J=\left(\begin{smallmatrix}0&-1\\1&0\end{smallmatrix}\right)\), this
means \(N=\left(\begin{smallmatrix}x&y\\y&-x\end{smallmatrix}\right)\),
so \(N^2=(x^2+y^2)I\) with \(x^2+y^2>0\).
\item From \(c_*^2=1\) we get \(M^2=\mu^2I\) with \(\mu>0\). Hence
either \(M\) is a scalar, and \(c\) acts trivially on \(D\), or \(M\)
has eigenvalues \(\mu,-\mu\), and \(\det M<0\).
\item On the other hand, \(c\) acts on \(D\) as \(\Ad(R^{k(c)})\). If
this action were nontrivial, we would have \(M=\lambda R^{k(c)}\) with
\(\lambda\in\R^\times\), so \(\det M=\lambda^2\,3^{-k(c)}>0\), a
contradiction.
\end{enumerate}
Hence \(c\) acts trivially on \(D\). As \(\Gamma\) acts faithfully,
\(c|_L=1\). So \(L\) is totally real, and in particular
\textbf{\(K\) is a real quadratic field}.

\subsection{Normalization}

We have \(\End^0_{\Q}(A)=D^\Gamma=\Q[R]\simeq\Q(\sqrt{-3})\). Put
\(u=3R-1\). Then \(u^2-u+1=0\), so \(u\) is a primitive sixth root of
unity, and
\[
R=\frac{1+u}{3},\qquad R^2=\frac u3,\qquad u^3=-1.
\]
Let \(\mathcal O=\Z[u]\simeq\Z[\zeta_3]\). Choose \(m\ge1\) with
\(\psi=mu\in\End_{\Q}(A)\), and let \(A'\subset A^2\) be the image of
\((m,\psi):A\to A^2\). It is \(\Q\)-isogenous to \(A\). The
endomorphism \(u'(x,y)=(y,y-x)\) of \(A^2\) satisfies
\(u'^2-u'+1=0\), and \(u'\circ(m,\psi)\circ[m]=(m,\psi)\circ\psi\),
because \(\psi^2=m\psi-m^2\). Hence \(u'(A')\subset A'\), and \(u'\)
corresponds to \(u\) under the isogeny \(A\to A'\). Thus
\(\mathcal O\subset\End_{\Q}(A')\). The surface \(A'\) is again of type
(S) with the same \(D\), \(L\) and \(\Gamma\), and
\(A'[p]^{\mathrm{ss}}\simeq A[p]^{\mathrm{ss}}\) for every prime \(p\).

\subsection{A twist with endomorphism field \texorpdfstring{\(K\)}{K}}

Define \(\xi(g)=u^{-2k(g)}\in\mu_3\subset\Aut_{\Q}(A')\). This is a
continuous homomorphism \(G_{\Q}\to\Aut_{\Q}(A')\), hence a
\(1\)-cocycle. Let \(B/\Q\) be the corresponding twist: there is an
isomorphism \(f:A'_{\overline\Q}\to B_{\overline\Q}\) with
\({}^gf=f\circ\xi(g)\) for all \(g\in G_{\Q}\).
\begin{enumerate}
\item Identify \(\End^0(B_{\overline\Q})\) with \(D\) by
\(\varphi\mapsto f\varphi f^{-1}\). The Galois action becomes
\[
\begin{gathered}
\varphi\longmapsto \xi(g)\,{}^g\varphi\,\xi(g)^{-1}
=\Ad\big(u^{-2k(g)}R^{k(g)}\big)\varphi=\Ad\big(S^{k(g)}\big)\varphi,\\
S=u^{-2}R,\qquad S^2=\frac{u^{-4}u}{3}=-\frac13 .
\end{gathered}
\]
\item Put \(\theta=3S=u^{-2}(1+u)\in\mathcal O\). Then
\(\theta^2=-3\).
\item The elements of \(\mathcal O\) commute with \(\xi(g)\) and with
\(R\), so they are fixed by the new action:
\(\mathcal O\subset\End_{\Q}(B)\). Since \(S^2\) is a scalar, all
geometric endomorphisms of \(B\) are defined over \(K\), and for
\(\tau\in G_{\Q}\setminus G_K\) we have
\({}^\tau\varphi=\theta\varphi\theta^{-1}\).
\item \(B_{\overline\Q}\simeq A'_{\overline\Q}\) is isogenous to the
square of a non-CM elliptic curve.
\item On Tate modules, \(f\) identifies \(T_\ell(B)\) with
\(T_\ell(A')\) on which \(g\) acts by \(x\mapsto\xi(g)\,g(x)\), since
\(g(f(x))={}^gf(g(x))=f(\xi(g)g(x))\).
\end{enumerate}

\subsection{A cyclic isogeny of degree 3}

\noindent\textbf{Lemma 8.4.}
\emph{Let \(\epsilon\in D\) be a rank-one idempotent and \(n\ge1\) an
integer with \(n\epsilon\in\End(B_{\overline\Q})\). Then
\(E_0=n\epsilon(B)\) is an elliptic curve over \(K\) without complex
multiplication. For \(\tau\in G_{\Q}\setminus G_K\) we have
\(\tau(E_0)=\theta(E_0)\ne E_0\), and
\(\theta|_{E_0}:E_0\to\tau(E_0)\) and
\(\theta|_{\tau(E_0)}:\tau(E_0)\to E_0\) are \(K\)-isogenies of degree
\(3\).}

\begin{proof}
The endomorphism \(n\epsilon\) is defined over \(K\), so \(E_0\) is an
elliptic curve over \(K\), isogenous to the non-CM curve \(E\). The
curve \(\tau(E_0)\) is the image of
\({}^\tau(n\epsilon)=\theta(n\epsilon)\theta^{-1}\), hence of
\(3\,{}^\tau(n\epsilon)=-\theta(n\epsilon)\theta\), whose image is
\(\theta(E_0)\). If \(\theta(E_0)=E_0\), then \(\theta|_{E_0}\) would
be an endomorphism of \(E_0\) with square \(-3\), and \(E_0\) would
have complex multiplication. Since \(\theta\) is defined over \(\Q\),
\(\theta|_{\tau E_0}={}^\tau(\theta|_{E_0})\), so both restrictions
have the same degree \(d\). Their composite is
\(\theta^2|_{E_0}=[-3]\), of degree \(9\). Hence \(d=3\).
\end{proof}

\section{The cyclic sextic case at 13}

\noindent\textbf{Proposition 9.1.}
\emph{Let \(A/\Q\) be of type (S) with \(\Gamma\simeq C_6\). Then
\(A[13]^{\mathrm{ss}}\) has no one-dimensional constituent. In
particular, \(A(\Q)[13]=0\).}

\begin{proof}
Keep the notation of Section~8, and suppose that
\(A[13]^{\mathrm{ss}}\) has a one-dimensional constituent.
\begin{enumerate}
\item \emph{The \(\mathfrak l\)-adic block.} Since
\(-3\equiv6^2\pmod{13}\), the prime \(13\) splits in
\(\Q(\sqrt{-3})\): \(\mathcal O\otimes\Z_{13}=\Z_{13}\times\Z_{13}\),
corresponding to primes \(\mathfrak l,\bar{\mathfrak l}\). The
idempotents of \(\mathcal O\otimes\Z_{13}\) lie in
\(\End_{\Q}(A')\otimes\Z_{13}\) and are rank-one idempotents of
\(D_{13}\), so
\(T_{13}(A')=T_{\mathfrak l}(A')\oplus T_{\bar{\mathfrak l}}(A')\) with
both summands \(G_{\Q}\)-stable of rank \(2\). Hence
\(A'[13]^{\mathrm{ss}}=M_{\mathfrak l}(A')^{\mathrm{ss}}\oplus
M_{\bar{\mathfrak l}}(A')^{\mathrm{ss}}\) with
\(M_{\mathfrak l}=T_{\mathfrak l}/13T_{\mathfrak l}\). Since
\(A'[13]^{\mathrm{ss}}\simeq A[13]^{\mathrm{ss}}\), after relabelling
\(M_{\mathfrak l}(A')\) has a one-dimensional constituent.
\item \emph{Transfer to \(B\).} Let \(T_{\mathfrak l}(B)\subset
T_{13}(B)\) be the image of the \(\mathfrak l\)-idempotent, which lies
in \(\End_{\Q}(B)\otimes\Z_{13}\). By Section~8.3,
\(T_{\mathfrak l}(B)\simeq T_{\mathfrak l}(A')\otimes\xi_{\mathfrak l}\),
where \(\xi_{\mathfrak l}:G_{\Q}\to\mu_3(\Z_{13})\) is the
\(\mathfrak l\)-component of \(\xi\). Hence
\(M_{\mathfrak l}(B)=T_{\mathfrak l}(B)/13\) is reducible: it contains
a \(G_{\Q}\)-stable line \(\overline{\mathcal C}\).
\item \emph{\(V_{13}(E_0)\) is the \(\mathfrak l\)-block.} Let
\(\iota:V_{13}(E_0)\to V_{13}(B)\to V_{\mathfrak l}(B)\) be inclusion
followed by the \(\mathfrak l\)-projection. It is \(G_K\)-equivariant.
By Faltings (Satz~3 and Satz~4), \(V_{13}(E_0)\) is a semisimple
\(G_K\)-representation with commutant
\(\End_K(E_0)\otimes\Q_{13}=\Q_{13}\), hence irreducible. So
\(\ker\iota\) is \(0\) or everything. In the second case
\(V_{13}(E_0)=V_{\bar{\mathfrak l}}(B)\) would be \(\theta\)-stable, so
\(V_{13}(\theta E_0)=\theta V_{13}(E_0)=V_{13}(E_0)\). But
\(\theta E_0\ne E_0\) by Lemma~8.4, so \(E_0+\theta E_0=B\) and
\(V_{13}(B)=V_{13}(E_0)\), which has the wrong dimension. Hence
\(\iota\) is an isomorphism.
\item \emph{The curve \(E_1\).} The lattice
\(\iota^{-1}(T_{\mathfrak l}(B))\) is \(G_K\)-stable. After scaling it
by a power of \(13\) it contains \(T_{13}(E_0)\), and its image in
\(V_{13}(E_0)/T_{13}(E_0)=E_0[13^\infty]\) is a finite \(G_K\)-stable
subgroup \(\Sigma\). Put \(E_1=E_0/\Sigma\) and let \(h:E_0\to E_1\) be
the quotient, a \(K\)-isogeny of degree \(13^b\) for some \(b\ge0\).
Then \(\iota\) and \(h\) give an isomorphism of \(G_K\)-modules
\(\iota_1:T_{13}(E_1)\to T_{\mathfrak l}(B)\). Put
\(\mathcal C=\iota_1^{-1}(\overline{\mathcal C})\subset E_1[13]\), a
\(G_K\)-stable subgroup of order \(13\). The curve \(E_1\) has no
complex multiplication.
\item \emph{An isogeny \(E_1^\tau\to E_1\) from Galois.} Fix
\(\tau\in G_{\Q}\setminus G_K\), and let \(\tau_B\) be the action of
\(\tau\) on \(T_{\mathfrak l}(B)\). Define
\[
\Phi_\tau=\iota_1^{-1}\circ\tau_B\circ\iota_1\circ\tau_*^{-1}:\
T_{13}(E_1^\tau)\longrightarrow T_{13}(E_1),
\]
where \(\tau_*:T_{13}(E_1)\to T_{13}(E_1^\tau)\) is induced by
\(\tau\). For \(g\in G_K\) we have \(\tau^{-1}g\tau\in G_K\), and
\(\iota_1\) is \(G_K\)-linear, so \(\Phi_\tau\) is
\(G_K\)-equivariant. The \(K\)-isogeny classes of \(E_1^\tau\) and
\(E_1\) coincide (both contain \(E_0\) and \(\tau E_0\)), and \(E_1\)
has no complex multiplication, so \(\Hom_K(E_1^\tau,E_1)=\Z\nu_0\) for
some isogeny \(\nu_0\). By Faltings (Korollar~1),
\(\Phi_\tau=\lambda\nu_0\) with \(\lambda\in\Z_{13}\). Since
\(\Phi_\tau\) is an isomorphism, \(\lambda\in\Z_{13}^\times\) and
\(13\nmid\deg\nu_0\).
\item \emph{\(\deg\nu_0=3\).} The \(K\)-isogeny
\[
h\circ\theta|_{\tau E_0}\circ\widehat{h^\tau}:\
E_1^\tau\longrightarrow\tau E_0\longrightarrow E_0\longrightarrow E_1
\]
has degree \(3\cdot13^{2b}\) and equals \(n\nu_0\) for some integer
\(n\). Hence \(n^2\deg\nu_0=3\cdot13^{2b}\). Since
\(13\nmid\deg\nu_0\), we get \(\deg\nu_0=3\).
\item \emph{\(\nu_0\) respects the level-\(13\) structure.} Since
\(\overline{\mathcal C}\) is \(G_{\Q}\)-stable,
\(\tau_B(\overline{\mathcal C})=\overline{\mathcal C}\), so
\(\Phi_\tau(\tau\mathcal C)=\mathcal C\) modulo \(13\). As \(\lambda\)
is a unit, \(\nu_0(\tau\mathcal C)=\mathcal C\).
\item \emph{A point of \(Y_0(39)(K)\).} Put
\(\nu_0^\tau={}^\tau\nu_0:E_1\to E_1^\tau\) (note
\(E_1^{\tau^2}=E_1\), since \(\tau^2\in G_K\)). The composite
\(\nu_0\circ\nu_0^\tau\) lies in \(\End_K(E_1)=\Z\) and has degree
\(9\), so it equals \(\pm3\). Hence
\(\nu_0^\tau(E_1[3])\subset\ker\nu_0\), and equality holds because both
groups have order \(3\) (the kernel of \(\nu_0^\tau\) has order \(3\)
and lies in \(E_1[3]\)). Applying \(\tau\) to
\(\nu_0(\tau\mathcal C)=\mathcal C\) gives
\(\nu_0^\tau(\mathcal C)=\tau\mathcal C\), because
\(\tau^2\mathcal C=\mathcal C\). Consider
\[
x=\big(E_1,\ \ker\nu_0^\tau+\mathcal C\big)\in Y_0(39)(K).
\]
Then \(\tau(x)=(E_1^\tau,\ker\nu_0+\tau\mathcal C)\), since
\(\tau(\ker\nu_0^\tau)=\ker({}^{\tau^2}\nu_0)=\ker\nu_0\). The
Atkin--Lehner involution \(w_3\) sends \((E,C_3+C_{13})\) to
\[
\big(E/C_3,\ E[3]/C_3+(C_{13}+C_3)/C_3\big).
\]
Using the isomorphism \(E_1/\ker\nu_0^\tau\simeq E_1^\tau\) induced by
\(\nu_0^\tau\), which maps \(E_1[3]\) onto \(\ker\nu_0\) and
\(\mathcal C\) onto \(\tau\mathcal C\), we obtain
\[
\tau(x)=w_3(x).
\]
\item \emph{\(x\) is not rational.} If \(x\in X_0(39)(\Q)\), then
\(x=\tau(x)=w_3(x)\). Then there is an isomorphism
\(\phi:E_1^\tau\to E_1\) over \(\overline\Q\), and
\(\phi\circ\nu_0^\tau\) is an endomorphism of \(E_1\) of degree \(3\),
which is impossible without complex multiplication. (Independently,
the computation below shows that the rational points of \(X_0(39)\)
are cusps.)
\item \emph{The computation.} Let \(\pi:X_0(39)\to X_0(39)/w_3\) be
the quotient map. Since \(\pi(\tau x)=\pi(w_3x)=\pi(x)\), the point
\(\pi(x)\) is fixed by \(\Gal(K/\Q)\), so it is rational. The script
\href{https://genus-2-torsion.github.io/data/knowledge/sources/geometric_26_39_x039_w3.m}{\texttt{geometric\_26\_39\_x039\_w3.m}}
(\href{https://genus-2-torsion.github.io/data/knowledge/sources/geometric_26_39_x039_w3.log}{log})
checks the following in Magma V2.30-1 \cite{Magma}, using the model of
\(X_0(39)\) in Magma's database of small modular curves:
\begin{itemize}
\item the model has genus \(3\), its \(j\)-map has degree \(56\) with
pole orders \(39,13,3,1\) at the four cusps, and
\(\#X(\F_p)=p+1-a_p\) for all primes \(p<100\), \(p\ne3,13\), where
\(a_p\) is the trace of the Hecke operator \(T_p\) on
\(S_2(\Gamma_0(39))\);
\item Magma's involution \(w_3\) is the moduli involution:
\(\Phi_3(j,j\circ w_3)=0\) in the function field, while
\(\Phi_{13}(j,j\circ w_3)\), \(\Phi_3(j,j\circ w_{13})\) and
\(\Phi_3(j,j\circ w_{39})\) are nonzero; here \(\Phi_n\) is the
classical modular polynomial. This identifies \(w_3\): the group \(\Gamma_0(39)\) permutes transitively the three subgroups of order \(3\) other than the one in the level structure, so \(j\circ w_3\) is the only root of \(\Phi_3(j,Y)\) in the function field of \(X_0(39)\); and an automorphism preserving \(j\) is trivial, since \(39\) is squarefree and so \(\Gamma_0(39)\) is its own normalizer in \(\mathrm{PSL}_2(\Z)\);
\item \(X_0(39)/w_3\) has genus \(1\) and is isomorphic over \(\Q\) to
the elliptic curve 39a1: \(y^2+xy=x^3+x^2-4x-5\). Its torsion is
\(\Z/2\times\Z/2\), its rational \(2\)-torsion is full and its
\(2\)-Selmer group has order \(4\), so its rank is \(0\). This is
unconditional: no class group of a number field enters the
\(2\)-descent. Hence \(X_0(39)/w_3\) has exactly four rational points;
\item over two of them, the fibre of \(\pi\) consists of two rational
cusps; over each of the other two, it consists of one pair of
conjugate points with residue field \(\Q(\sqrt{-7})\). In particular
\(X_0(39)(\Q)\) consists of the four cusps.
\end{itemize}
The point \(x\) is not rational and lies over a rational point of
\(X_0(39)/w_3\). Hence its residue field is \(\Q(\sqrt{-7})\), and
\(\Q(\sqrt{-7})\subset K\). So \(K=\Q(\sqrt{-7})\), contradicting
Section~8.1, where we showed that \(K\) is real.\qedhere
\end{enumerate}
\end{proof}

\section{The cyclic sextic case at 3}

\noindent\textbf{Lemma 10.1.}
\emph{Let \(A/\Q\) be of type (S) with \(\Gamma\simeq C_6\). Then
\(A[3]^{\mathrm{ss}}\) has no trivial constituent. In particular,
\(A(\Q)[3]=0\).}

\begin{proof}
Keep the notation of Section~8, and suppose that \(A[3]^{\mathrm{ss}}\)
has a trivial constituent.
\begin{enumerate}
\item \emph{Reduction to \(A'[\theta]\).} The ring
\(\mathcal O_3=\mathcal O\otimes\Z_3=\Z_3[\zeta_3]\) is a discrete
valuation ring with uniformizer \(\theta\), and \(T_3(A')\) is free of
rank \(2\) over it. In the filtration
\[
\theta T_3(A')/3T_3(A')\ \subset\ T_3(A')/3T_3(A')=A'[3],
\]
both graded pieces are isomorphic to
\(T_3(A')/\theta T_3(A')\simeq A'[\theta]\), via the \(G_{\Q}\)-linear
map \(\theta\). Hence
\(A'[3]^{\mathrm{ss}}\simeq(A'[\theta]^{\mathrm{ss}})^{\oplus2}\), and
\(A'[\theta]^{\mathrm{ss}}\) has a trivial constituent.
\item \emph{Passing to \(B[\theta]\).} Since
\(\zeta_3\equiv1\pmod{\theta}\), each \(\xi(g)\in\mu_3\) acts trivially
on \(A'[\theta]\). By Section~8.3, \(B[\theta]\simeq A'[\theta]\) as
\(G_{\Q}\)-modules, so \(B[\theta]^{\mathrm{ss}}\) has a trivial
constituent.
\item \emph{Two lines in \(B[\theta]\).} The group \(B[\theta]\) has
order \(\deg\theta=9\) and is killed by \(\theta^2=-3\), so it is a
\(2\)-dimensional \(\F_3\)-vector space. Put
\(\Delta=\ker(\theta|_{E_0})\), a \(G_K\)-stable line in \(B[\theta]\),
on which \(G_K\) acts by a character \(\alpha\). Fix
\(\tau\in G_{\Q}\setminus G_K\). Then
\(\tau\Delta=\ker(\theta|_{\tau E_0})\), and \(G_K\) acts on it by
\(\alpha^\tau(g)=\alpha(\tau^{-1}g\tau)\). On the other hand
\(\ker(\theta|_{\tau E_0})=\theta(E_0[3])\): this image lies in
\(\tau E_0=\theta E_0\), it is killed by \(\theta|_{\tau E_0}\) because
\(\theta^2=-3\), and both groups have order \(3\). As \(G_K\)-modules,
\(\theta(E_0[3])\simeq E_0[3]/\Delta\), and
\(E_0[3]/\Delta\simeq\Hom(\Delta,\mu_3)\) by the Weil pairing. Hence
\[
\alpha^\tau=\alpha^{-1}\,\bar\chi_3|_{G_K}.
\]
\item \emph{Case \(\tau\Delta=\Delta\).} Then \(\alpha^\tau=\alpha\),
so \(\bar\chi_3|_{G_K}=\alpha^2=1\), because \(\alpha\) takes values in
\(\F_3^\times=\{\pm1\}\).
\item \emph{Case \(\tau\Delta\ne\Delta\).} Then
\(B[\theta]=\Delta\oplus\tau\Delta\), with \(G_K\)-composition factors
\(\alpha\) and \(\alpha^\tau\). A trivial \(G_{\Q}\)-composition factor
restricts to a trivial \(G_K\)-composition factor, so \(\alpha=1\) or
\(\alpha^\tau=1\). Each of these implies the other, so
\(\bar\chi_3|_{G_K}=\alpha\alpha^\tau=1\).
\end{enumerate}
In both cases \(\zeta_3\in K\), that is, \(K=\Q(\sqrt{-3})\). This
contradicts Section~8.1.
\end{proof}

\section{Completion of the proof and remarks}

Let \(P\in A(\Q)\) have order \(N\in\{26,39\}\). Type (Q) and type
(N) are impossible by Section~3, and type (CM) by Proposition~4.2. If
\(A\) has type (S), Proposition~7.1 gives \(\Gamma\simeq C_6\), and
then Proposition~9.1, applied to \(P_{13}\), gives a contradiction. For
\(N=39\), Lemma~10.1, applied to \(P_3\), gives a second
contradiction, which does not use the computation on \(X_0(39)\). This
proves Theorem~\ref{thm:main}.\hfill\(\square\)

\medskip\noindent\textbf{Remarks.}
\begin{itemize}
\item The point of order \(2\) or \(3\) is used only in type (N) and
for \(\Gamma\in\{C_2,C_3,C_4\}\), through Lemma~2.1, and in type (CM),
where the whole point of order \(N\) enters the point count of
Proposition~4.2. The argument
excludes rational \(13\)-torsion by itself only in type (Q), for
\(\Gamma\in\{D_4,D_6,D_8,D_{12}\}\), and for \(\Gamma=C_6\).
\item Torsion is never moved along an arbitrary isogeny. Only residual
semisimplifications are compared, and points are combined only inside
one elliptic isogeny class over one field (Lemma~2.1).
\item The census curve 19.b, \(y^2+(x^3+x^2+1)y=x^2+x\), has Jacobian
\(J_1(13)\) (LMFDB label 169.a.169.1 \cite{LMFDB}). The LMFDB lists its
Sato--Tate group as \(E_6\) and its endomorphism field as the real
cyclic sextic field \(\Q(\zeta_{13})^+\), so it has type (S) with
\(\Gamma\simeq C_6\), and \(J_1(13)(\Q)\simeq\Z/19\). The supplementary
script checks that the characteristic polynomial of Frobenius at \(2\)
has no root modulo \(13\) and that \(\#J_1(13)(\F_2)=19\). Hence
\(J_1(13)[13]^{\mathrm{ss}}\) has no one-dimensional constituent and
\(J_1(13)[3]^{\mathrm{ss}}\) has no trivial constituent, as
Proposition~9.1 and Lemma~10.1 require.
\item The fibres of \(X_0(39)\to X_0(39)/w_3\) above the two
non-cuspidal rational points consist of non-CM points over the
imaginary field \(\Q(\sqrt{-7})\) that \(w_3\) maps to their
conjugates. Thus the reality of \(K\) (Section~8.1) is essential in
Proposition~9.1.
\end{itemize}

\section{Consequence for the census}

The groups \href{https://genus-2-torsion.github.io/group.html?g=26}{\([26]\)},
\href{https://genus-2-torsion.github.io/group.html?g=2.26}{\([2,26]\)} and
\href{https://genus-2-torsion.github.io/group.html?g=39}{\([39]\)} contain an element of order
\(26\) or \(39\). They are therefore impossible in the class
``geometrically split, simple over \(\Q\)''. They are also impossible
in the class ``split over \(\Q\)'' by the
\href{https://genus-2-torsion.github.io/proofs/qsplit-prime-obstruction.html}{prime torsion
obstruction}, since \(13\) divides their orders. Hence their combined
geometrically split \(\infty\)? cells are \(\emptyset\). The argument
applies to every geometric splitting degree, to every polarization, and
to every integral endomorphism ring.

The proof of
\href{https://genus-2-torsion.github.io/proofs/qsimple-13-obstruction.html}{the \([13]\)
obstruction}, which treats rational \(13\)-torsion on its own, gives an
independent argument for these three groups in the class
``geometrically split, simple over \(\Q\)''. The present proof does not
use it.

\begin{thebibliography}{99}
\bibitem{Mazur}
B. Mazur,
\emph{Modular curves and the Eisenstein ideal},
Publ. Math. IH\'ES \textbf{47} (1977), 33--186, Theorem~(7$'$) (p.~35).
\href{https://www.numdam.org/item/PMIHES_1977__47__33_0/}{Numdam}.

\bibitem{KM}
M. A. Kenku and F. Momose,
\emph{Torsion points on elliptic curves defined over quadratic fields},
Nagoya Math. J. \textbf{109} (1988), 125--149, Theorem~(0.1).
\href{https://doi.org/10.1017/S0027763000002816}{doi:10.1017/S0027763000002816}.

\bibitem{DEvHMZB}
M. Derickx, A. Etropolski, M. van Hoeij, J. S. Morrow and
D. Zureick-Brown,
\emph{Sporadic cubic torsion},
Algebra \& Number Theory \textbf{15} (2021), 1837--1864, Theorem~A.
\href{https://doi.org/10.2140/ant.2021.15.1837}{doi:10.2140/ant.2021.15.1837};
\href{https://arxiv.org/abs/2007.13929}{arXiv:2007.13929}.

\bibitem{DN}
M. Derickx and F. Najman,
\emph{Classification of torsion of elliptic curves over quartic fields},
arXiv:2412.16016v2 (2025), Theorem~1.1.
\href{https://arxiv.org/pdf/2412.16016}{arXiv:2412.16016}.

\bibitem{Faltings}
G. Faltings,
\emph{Endlichkeitss\"atze f\"ur abelsche Variet\"aten \"uber
Zahlk\"orpern},
Invent. Math. \textbf{73} (1983), 349--366, \S5, Satz~3, Satz~4 and
Korollar~1 (pp.~360--361).
\href{https://doi.org/10.1007/BF01388432}{doi:10.1007/BF01388432};
Erratum, Invent. Math. \textbf{75} (1984), 381,
\href{https://doi.org/10.1007/BF01388572}{doi:10.1007/BF01388572}.

\bibitem{SerreTate}
J.-P. Serre and J. Tate,
\emph{Good reduction of abelian varieties},
Ann. of Math. \textbf{88} (1968), 492--517, \S1, Theorem~1 (p.~493), \S2, p.~498, and \S5,
Theorem~6(a) (p.~503).
\href{https://doi.org/10.2307/1970722}{doi:10.2307/1970722}.

\bibitem{Clark}
P. L. Clark,
\emph{Bounds for torsion on abelian varieties with integral moduli},
arXiv:math/0407264v2 (2004), \S2.1 and Theorem~5.
\href{https://arxiv.org/pdf/math/0407264}{arXiv:math/0407264}.

\bibitem{ClarkXarles}
P. L. Clark and X. Xarles,
\emph{Local bounds for torsion points on abelian varieties},
Canad. J. Math. \textbf{60} (2008), 532--555: Main Theorem~(i) and
\S2.4 for the local bound, and Theorem~1.4 for the list, stated there
for abelian surfaces with everywhere anisotropic reduction, which
includes potentially good reduction.
\href{https://doi.org/10.4153/CJM-2008-026-x}{doi:10.4153/CJM-2008-026-x}.

\bibitem{Magma}
W. Bosma, J. Cannon and C. Playoust,
\emph{The Magma algebra system I: the user language},
J. Symbolic Comput. \textbf{24} (1997), 235--265.
\href{https://doi.org/10.1006/jsco.1996.0125}{doi:10.1006/jsco.1996.0125}.
Computations in Magma V2.30-1, including its database of small modular
curves.

\bibitem{LMFDB}
The LMFDB Collaboration,
\emph{The L-functions and modular forms database}, genus 2 curve
169.a.169.1 (used only in the remarks of Section~11).
\href{https://www.lmfdb.org/Genus2Curve/Q/169/a/169/1}{LMFDB}.
\end{thebibliography}

\end{document}
